% --- Archon physics formalization source begin ---
% archon:physics
% archon:covers PhyXMiniProblems/problem_phyx_mini_0852.lean
% archon:source-report reports/phyx_mini/problem_phyx_mini_0852.source.json

\chapter{Physics problem phyx\_mini\_0852}
\label{ch:PhyXMiniProblems_problem_phyx_mini_0852}

\paragraph{Problem source.}
\#\# Physical scenario

The figure shows a solid metal sphere at the center of a hollow metal sphere.

\#\# Question

What is the total charge on the inside surface of the hollow sphere?

\#\# Answer choices

- A: A: 80C

- B: B: 10C

- C: C: 20C

- D: D: 0C

\#\# Figure caption (auxiliary; use the image as primary evidence)

The image depicts two concentric conducting spheres with specified radii and labeled distances. The inner sphere has a radius of 5 cm, while the outer sphere has a radius of 10 cm. The outer boundary is delineated at a radius of 15 cm. 

There are two electric field vectors (\textbackslash{}( \textbackslash{}vec\{E\} \textbackslash{})) shown, each with a magnitude of 15,000 N/C. One vector points inward towards the center from outside the outer sphere, and the other points outward from the inner to the outer sphere.

The distances from the center are marked as follows:
- From the center to the inner sphere: 5 cm
- From the center to the outer sphere: 10 cm
- The outermost boundary extends to 15 cm.

Additionally, there are horizontal measurements:
- A distance labeled as 8 cm within the boundary of the outer sphere
- The total distance from the center to a point outside the outer field measurement is 17 cm.

The term "Conducting spheres" is annotated pointing to the spheres, indicating their material.

\#\# Dataset metadata

Electrostatics; Physical Model Grounding Reasoning, Multi-Formula Reasoning

\paragraph{Recorded answer/context.}
D: D: 0C

\paragraph{Figure/image path.}
/root/proposal\_for\_physic/science-mango/phyx\_mini\_run/phyx\_data/test\_image/852.png

\paragraph{Formalization target.}
create a compiling Lean file with sorry bodies at `PhyXMiniProblems/problem\_phyx\_mini\_0852.lean`. The Lean declarations must preserve the physical quantities, dimensions or dimensional roles, figure labels, governing-law hypotheses, and final relation expressed by this problem.
Use Mathlib/Physlib names found through LeanExplore where available. If a physics API is missing, introduce faithful local abstractions rather than scalar placeholder aliases.

\begin{theorem}[Physics formalization target]
\label{thm:physics:phyx_mini_0852:target}
\lean{PhyXMiniProblems.ProblemPhyXMini0852.problem_phyx_mini_0852}
\uses{def:physics:phyx-mini-0852:phyxminiproblems-problemphyxmini0852-chargeincoulombs, def:physics:phyx-mini-0852:phyxminiproblems-problemphyxmini0852-electricpermittivityinsi, def:physics:phyx-mini-0852:phyxminiproblems-problemphyxmini0852-concentricconductingspheressetup, def:physics:phyx-mini-0852:phyxminiproblems-problemphyxmini0852-matchesconcentricconductingspherescenario, def:physics:phyx-mini-0852:phyxminiproblems-problemphyxmini0852-matchesprimaryconcentricspheresfigure, def:physics:phyx-mini-0852:phyxminiproblems-problemphyxmini0852-satisfiesconcentricsphericalgeometry, def:physics:phyx-mini-0852:phyxminiproblems-problemphyxmini0852-fieldsamplescalibrateelectricfield, def:physics:phyx-mini-0852:phyxminiproblems-problemphyxmini0852-satisfieselectrostaticconductorequilibrium, def:physics:phyx-mini-0852:phyxminiproblems-problemphyxmini0852-satisfiessphericalgausslaw, def:physics:phyx-mini-0852:phyxminiproblems-problemphyxmini0852-hasphysicalelectrostaticparameters, def:physics:phyx-mini-0852:phyxminiproblems-problemphyxmini0852-answerchargeincoulombs, def:physics:phyx-mini-0852:phyxminiproblems-problemphyxmini0852-recordeddatasetanswer}
The assigned autoformalize agent should translate this physics problem into Lean declarations in the covered file, with theorem and lemma proof bodies written as `by sorry`.
\end{theorem}
\begin{proof}
This is an autoformalization task, not a proof task. Produce statements that can later be proved by the physics prover without weakening the physical model.
\end{proof}

% --- Archon declaration topology begin ---
\section{Declaration topology}
\begin{definition}[electric Field Strength Dimension]
  \label{def:physics:phyx-mini-0852:phyxminiproblems-problemphyxmini0852-electricfieldstrengthdimension}
  \lean{PhyXMiniProblems.ProblemPhyXMini0852.electricFieldStrengthDimension}
  The physical dimension M L T\textasciicircum{}-2 C\textasciicircum{}-1 of electric-field strength
\end{definition}
\begin{proof}
  This is the declared model definition of electric Field Strength Dimension.
\end{proof}
\begin{definition}[electric Permittivity Dimension]
  \label{def:physics:phyx-mini-0852:phyxminiproblems-problemphyxmini0852-electricpermittivitydimension}
  \lean{PhyXMiniProblems.ProblemPhyXMini0852.electricPermittivityDimension}
  The physical dimension C\textasciicircum{}2 T\textasciicircum{}2 M\textasciicircum{}-1 L\textasciicircum{}-3 of permittivity
\end{definition}
\begin{proof}
  This is the declared model definition of electric Permittivity Dimension.
\end{proof}
\begin{definition}[Length Quantity]
  \label{def:physics:phyx-mini-0852:phyxminiproblems-problemphyxmini0852-lengthquantity}
  \lean{PhyXMiniProblems.ProblemPhyXMini0852.LengthQuantity}
  A nonnegative, unit-independent physical length
\end{definition}
\begin{proof}
  This is the declared model definition of Length Quantity.
\end{proof}
\begin{definition}[Signed Charge Quantity]
  \label{def:physics:phyx-mini-0852:phyxminiproblems-problemphyxmini0852-signedchargequantity}
  \lean{PhyXMiniProblems.ProblemPhyXMini0852.SignedChargeQuantity}
  A signed, unit-independent physical electric charge
\end{definition}
\begin{proof}
  This is the declared model definition of Signed Charge Quantity.
\end{proof}
\begin{definition}[Electric Field Magnitude Quantity]
  \label{def:physics:phyx-mini-0852:phyxminiproblems-problemphyxmini0852-electricfieldmagnitudequantity}
  \lean{PhyXMiniProblems.ProblemPhyXMini0852.ElectricFieldMagnitudeQuantity}
  \uses{def:physics:phyx-mini-0852:phyxminiproblems-problemphyxmini0852-electricfieldstrengthdimension}
  A nonnegative, unit-independent electric-field magnitude
\end{definition}
\begin{proof}
  This is the declared model definition of Electric Field Magnitude Quantity.
\end{proof}
\begin{definition}[Electric Permittivity Quantity]
  \label{def:physics:phyx-mini-0852:phyxminiproblems-problemphyxmini0852-electricpermittivityquantity}
  \lean{PhyXMiniProblems.ProblemPhyXMini0852.ElectricPermittivityQuantity}
  \uses{def:physics:phyx-mini-0852:phyxminiproblems-problemphyxmini0852-electricpermittivitydimension}
  A nonnegative, unit-independent vacuum permittivity
\end{definition}
\begin{proof}
  This is the declared model definition of Electric Permittivity Quantity.
\end{proof}
\begin{definition}[length In Meters]
  \label{def:physics:phyx-mini-0852:phyxminiproblems-problemphyxmini0852-lengthinmeters}
  \lean{PhyXMiniProblems.ProblemPhyXMini0852.lengthInMeters}
  \uses{def:physics:phyx-mini-0852:phyxminiproblems-problemphyxmini0852-lengthquantity}
  Coherent-SI readout of a physical length, in metres
\end{definition}
\begin{proof}
  This is the declared model definition of length In Meters.
\end{proof}
\begin{definition}[length In Centimeters]
  \label{def:physics:phyx-mini-0852:phyxminiproblems-problemphyxmini0852-lengthincentimeters}
  \lean{PhyXMiniProblems.ProblemPhyXMini0852.lengthInCentimeters}
  \uses{def:physics:phyx-mini-0852:phyxminiproblems-problemphyxmini0852-lengthquantity, def:physics:phyx-mini-0852:phyxminiproblems-problemphyxmini0852-lengthinmeters}
  Centimetre readout used by all five distance labels in the raster
\end{definition}
\begin{proof}
  This is the declared model definition of length In Centimeters.
\end{proof}
\begin{definition}[charge In Coulombs]
  \label{def:physics:phyx-mini-0852:phyxminiproblems-problemphyxmini0852-chargeincoulombs}
  \lean{PhyXMiniProblems.ProblemPhyXMini0852.chargeInCoulombs}
  \uses{def:physics:phyx-mini-0852:phyxminiproblems-problemphyxmini0852-signedchargequantity}
  Coherent-SI readout of a signed electric charge, in coulombs
\end{definition}
\begin{proof}
  This is the declared model definition of charge In Coulombs.
\end{proof}
\begin{definition}[field Magnitude In Newtons Per Coulomb]
  \label{def:physics:phyx-mini-0852:phyxminiproblems-problemphyxmini0852-fieldmagnitudeinnewtonspercoulomb}
  \lean{PhyXMiniProblems.ProblemPhyXMini0852.fieldMagnitudeInNewtonsPerCoulomb}
  \uses{def:physics:phyx-mini-0852:phyxminiproblems-problemphyxmini0852-electricfieldmagnitudequantity}
  Coherent-SI readout of field magnitude, in newtons per coulomb
\end{definition}
\begin{proof}
  This is the declared model definition of field Magnitude In Newtons Per Coulomb.
\end{proof}
\begin{definition}[electric Permittivity In SI]
  \label{def:physics:phyx-mini-0852:phyxminiproblems-problemphyxmini0852-electricpermittivityinsi}
  \lean{PhyXMiniProblems.ProblemPhyXMini0852.electricPermittivityInSI}
  \uses{def:physics:phyx-mini-0852:phyxminiproblems-problemphyxmini0852-electricpermittivityquantity}
  Coherent-SI readout of electric permittivity
\end{definition}
\begin{proof}
  This is the declared model definition of electric Permittivity In SI.
\end{proof}
\begin{definition}[Sphere Boundary]
  \label{def:physics:phyx-mini-0852:phyxminiproblems-problemphyxmini0852-sphereboundary}
  \lean{PhyXMiniProblems.ProblemPhyXMini0852.SphereBoundary}
  The three spherical boundaries explicitly labelled in the image
\end{definition}
\begin{proof}
  This is the declared model definition of Sphere Boundary.
\end{proof}
\begin{definition}[Radial Field Sample]
  \label{def:physics:phyx-mini-0852:phyxminiproblems-problemphyxmini0852-radialfieldsample}
  \lean{PhyXMiniProblems.ProblemPhyXMini0852.RadialFieldSample}
  The two black-dot locations at which electric-field arrows are drawn
\end{definition}
\begin{proof}
  This is the declared model definition of Radial Field Sample.
\end{proof}
\begin{definition}[Radial Arrow Orientation]
  \label{def:physics:phyx-mini-0852:phyxminiproblems-problemphyxmini0852-radialarroworientation}
  \lean{PhyXMiniProblems.ProblemPhyXMini0852.RadialArrowOrientation}
  Orientation of a field arrow relative to the radius through its point
\end{definition}
\begin{proof}
  This is the declared model definition of Radial Arrow Orientation.
\end{proof}
\begin{definition}[Concentric Conducting Spheres Figure]
  \label{def:physics:phyx-mini-0852:phyxminiproblems-problemphyxmini0852-concentricconductingspheresfigure}
  \lean{PhyXMiniProblems.ProblemPhyXMini0852.ConcentricConductingSpheresFigure}
  \uses{def:physics:phyx-mini-0852:phyxminiproblems-problemphyxmini0852-lengthquantity, def:physics:phyx-mini-0852:phyxminiproblems-problemphyxmini0852-electricfieldmagnitudequantity, def:physics:phyx-mini-0852:phyxminiproblems-problemphyxmini0852-sphereboundary, def:physics:phyx-mini-0852:phyxminiproblems-problemphyxmini0852-radialfieldsample, def:physics:phyx-mini-0852:phyxminiproblems-problemphyxmini0852-radialarroworientation}
  Typed content of the supplied raster. Distance and field labels are physical quantities; their printed numerical readouts are stated separately below
\end{definition}
\begin{proof}
  This is the declared model definition of Concentric Conducting Spheres Figure.
\end{proof}
\begin{definition}[Concentric Conducting Spheres Setup]
  \label{def:physics:phyx-mini-0852:phyxminiproblems-problemphyxmini0852-concentricconductingspheressetup}
  \lean{PhyXMiniProblems.ProblemPhyXMini0852.ConcentricConductingSpheresSetup}
  \uses{def:physics:phyx-mini-0852:phyxminiproblems-problemphyxmini0852-lengthquantity, def:physics:phyx-mini-0852:phyxminiproblems-problemphyxmini0852-signedchargequantity, def:physics:phyx-mini-0852:phyxminiproblems-problemphyxmini0852-electricfieldmagnitudequantity, def:physics:phyx-mini-0852:phyxminiproblems-problemphyxmini0852-electricpermittivityquantity, def:physics:phyx-mini-0852:phyxminiproblems-problemphyxmini0852-radialfieldsample, def:physics:phyx-mini-0852:phyxminiproblems-problemphyxmini0852-concentricconductingspheresfigure}
  Independent physical objects in the concentric-sphere experiment. In particular, hollowInnerSurfaceCharge is an unconstrained physical field: it is not defined from answer D or from zero
\end{definition}
\begin{proof}
  This is the declared model definition of Concentric Conducting Spheres Setup.
\end{proof}
\begin{definition}[sample Radius]
  \label{def:physics:phyx-mini-0852:phyxminiproblems-problemphyxmini0852-sampleradius}
  \lean{PhyXMiniProblems.ProblemPhyXMini0852.sampleRadius}
  \uses{def:physics:phyx-mini-0852:phyxminiproblems-problemphyxmini0852-lengthquantity, def:physics:phyx-mini-0852:phyxminiproblems-problemphyxmini0852-radialfieldsample, def:physics:phyx-mini-0852:phyxminiproblems-problemphyxmini0852-concentricconductingspheressetup}
  The physical radius associated with either field-sampling point
\end{definition}
\begin{proof}
  This is the declared model definition of sample Radius.
\end{proof}
\begin{definition}[spherical Flux In Newton Square Meters Per Coulomb]
  \label{def:physics:phyx-mini-0852:phyxminiproblems-problemphyxmini0852-sphericalfluxinnewtonsquaremeterspercoulomb}
  \lean{PhyXMiniProblems.ProblemPhyXMini0852.sphericalFluxInNewtonSquareMetersPerCoulomb}
  \uses{def:physics:phyx-mini-0852:phyxminiproblems-problemphyxmini0852-lengthinmeters, def:physics:phyx-mini-0852:phyxminiproblems-problemphyxmini0852-fieldmagnitudeinnewtonspercoulomb, def:physics:phyx-mini-0852:phyxminiproblems-problemphyxmini0852-radialfieldsample, def:physics:phyx-mini-0852:phyxminiproblems-problemphyxmini0852-concentricconductingspheressetup, def:physics:phyx-mini-0852:phyxminiproblems-problemphyxmini0852-sampleradius}
  For a radial field of magnitude E on a sphere of radius r, the outward electric flux is 4 π r² E. This helper is only the flux expression; Gauss's law is imposed independently in SatisfiesSphericalGaussLaw
\end{definition}
\begin{proof}
  This is the declared model definition of spherical Flux In Newton Square Meters Per Coulomb.
\end{proof}
\begin{definition}[Matches Concentric Conducting Sphere Scenario]
  \label{def:physics:phyx-mini-0852:phyxminiproblems-problemphyxmini0852-matchesconcentricconductingspherescenario}
  \lean{PhyXMiniProblems.ProblemPhyXMini0852.MatchesConcentricConductingSphereScenario}
  \uses{def:physics:phyx-mini-0852:phyxminiproblems-problemphyxmini0852-concentricconductingspheressetup}
  The qualitative material and nesting claims made by the problem prose
\end{definition}
\begin{proof}
  This is the declared model definition of Matches Concentric Conducting Sphere Scenario.
\end{proof}
\begin{definition}[Matches Primary Concentric Spheres Figure]
  \label{def:physics:phyx-mini-0852:phyxminiproblems-problemphyxmini0852-matchesprimaryconcentricspheresfigure}
  \lean{PhyXMiniProblems.ProblemPhyXMini0852.MatchesPrimaryConcentricSpheresFigure}
  \uses{def:physics:phyx-mini-0852:phyxminiproblems-problemphyxmini0852-lengthincentimeters, def:physics:phyx-mini-0852:phyxminiproblems-problemphyxmini0852-fieldmagnitudeinnewtonspercoulomb, def:physics:phyx-mini-0852:phyxminiproblems-problemphyxmini0852-concentricconductingspheressetup}
  Primary-raster evidence, including both radially outward arrows. Each printed label is first identified with an independent physical quantity and is only then assigned its centimetre or N/C readout
\end{definition}
\begin{proof}
  This is the declared model definition of Matches Primary Concentric Spheres Figure.
\end{proof}
\begin{definition}[Satisfies Concentric Spherical Geometry]
  \label{def:physics:phyx-mini-0852:phyxminiproblems-problemphyxmini0852-satisfiesconcentricsphericalgeometry}
  \lean{PhyXMiniProblems.ProblemPhyXMini0852.SatisfiesConcentricSphericalGeometry}
  \uses{def:physics:phyx-mini-0852:phyxminiproblems-problemphyxmini0852-lengthinmeters, def:physics:phyx-mini-0852:phyxminiproblems-problemphyxmini0852-concentricconductingspheressetup, def:physics:phyx-mini-0852:phyxminiproblems-problemphyxmini0852-sampleradius}
  The concentric radial geometry encoded by the five labelled radii. The conducting regions are described by distance from the common center, chosen as the origin of Space 3
\end{definition}
\begin{proof}
  This is the declared model definition of Satisfies Concentric Spherical Geometry.
\end{proof}
\begin{definition}[Field Samples Calibrate Electric Field]
  \label{def:physics:phyx-mini-0852:phyxminiproblems-problemphyxmini0852-fieldsamplescalibrateelectricfield}
  \lean{PhyXMiniProblems.ProblemPhyXMini0852.FieldSamplesCalibrateElectricField}
  \uses{def:physics:phyx-mini-0852:phyxminiproblems-problemphyxmini0852-fieldmagnitudeinnewtonspercoulomb, def:physics:phyx-mini-0852:phyxminiproblems-problemphyxmini0852-concentricconductingspheressetup}
  The dimensionful field magnitudes calibrate the Physlib vector field at the two black-dot samples. The positive scalar multiple records radial outward orientation independently of the pictorial arrow fields
\end{definition}
\begin{proof}
  This is the declared model definition of Field Samples Calibrate Electric Field.
\end{proof}
\begin{definition}[Satisfies Electrostatic Conductor Equilibrium]
  \label{def:physics:phyx-mini-0852:phyxminiproblems-problemphyxmini0852-satisfieselectrostaticconductorequilibrium}
  \lean{PhyXMiniProblems.ProblemPhyXMini0852.SatisfiesElectrostaticConductorEquilibrium}
  \uses{def:physics:phyx-mini-0852:phyxminiproblems-problemphyxmini0852-chargeincoulombs, def:physics:phyx-mini-0852:phyxminiproblems-problemphyxmini0852-concentricconductingspheressetup}
  Electrostatic equilibrium of both conductors. The total charge of the hollow sphere is decomposed into independent inner- and outer-surface charges; no surface charge is assigned the requested numerical answer here
\end{definition}
\begin{proof}
  This is the declared model definition of Satisfies Electrostatic Conductor Equilibrium.
\end{proof}
\begin{definition}[Satisfies Spherical Gauss Law]
  \label{def:physics:phyx-mini-0852:phyxminiproblems-problemphyxmini0852-satisfiessphericalgausslaw}
  \lean{PhyXMiniProblems.ProblemPhyXMini0852.SatisfiesSphericalGaussLaw}
  \uses{def:physics:phyx-mini-0852:phyxminiproblems-problemphyxmini0852-chargeincoulombs, def:physics:phyx-mini-0852:phyxminiproblems-problemphyxmini0852-electricpermittivityinsi, def:physics:phyx-mini-0852:phyxminiproblems-problemphyxmini0852-concentricconductingspheressetup, def:physics:phyx-mini-0852:phyxminiproblems-problemphyxmini0852-sphericalfluxinnewtonsquaremeterspercoulomb}
  Spherical Gauss-law balances in coherent SI units. The middle equation is the zero-flux Gaussian surface lying in the metal of the shell. It states the standard induced-charge cancellation law, not the requested 0 C answer. The exterior balance retains all three surface contributions
\end{definition}
\begin{proof}
  This is the declared model definition of Satisfies Spherical Gauss Law.
\end{proof}
\begin{definition}[Has Physical Electrostatic Parameters]
  \label{def:physics:phyx-mini-0852:phyxminiproblems-problemphyxmini0852-hasphysicalelectrostaticparameters}
  \lean{PhyXMiniProblems.ProblemPhyXMini0852.HasPhysicalElectrostaticParameters}
  \uses{def:physics:phyx-mini-0852:phyxminiproblems-problemphyxmini0852-electricpermittivityinsi, def:physics:phyx-mini-0852:phyxminiproblems-problemphyxmini0852-concentricconductingspheressetup}
  Positivity of the vacuum permittivity used by the Gauss-law balances
\end{definition}
\begin{proof}
  This is the declared model definition of Has Physical Electrostatic Parameters.
\end{proof}
\begin{definition}[Answer Choice]
  \label{def:physics:phyx-mini-0852:phyxminiproblems-problemphyxmini0852-answerchoice}
  \lean{PhyXMiniProblems.ProblemPhyXMini0852.AnswerChoice}
  Labels of the four charge choices printed in the source record
\end{definition}
\begin{proof}
  This is the declared model definition of Answer Choice.
\end{proof}
\begin{definition}[answer Charge In Coulombs]
  \label{def:physics:phyx-mini-0852:phyxminiproblems-problemphyxmini0852-answerchargeincoulombs}
  \lean{PhyXMiniProblems.ProblemPhyXMini0852.answerChargeInCoulombs}
  \uses{def:physics:phyx-mini-0852:phyxminiproblems-problemphyxmini0852-answerchoice}
  Coulomb number printed beside each answer choice
\end{definition}
\begin{proof}
  This is the declared model definition of answer Charge In Coulombs.
\end{proof}
\begin{definition}[recorded Dataset Answer]
  \label{def:physics:phyx-mini-0852:phyxminiproblems-problemphyxmini0852-recordeddatasetanswer}
  \lean{PhyXMiniProblems.ProblemPhyXMini0852.recordedDatasetAnswer}
  \uses{def:physics:phyx-mini-0852:phyxminiproblems-problemphyxmini0852-answerchoice}
  The answer label recorded by the dataset
\end{definition}
\begin{proof}
  This is the declared model definition of recorded Dataset Answer.
\end{proof}
\begin{lemma}[hollow inner surface charge cancels central]
  \label{lem:physics:phyx-mini-0852:phyxminiproblems-problemphyxmini0852-hollow-inner-surface-charge-cancels-central}
  \lean{PhyXMiniProblems.ProblemPhyXMini0852.hollow_inner_surface_charge_cancels_central}
  \uses{def:physics:phyx-mini-0852:phyxminiproblems-problemphyxmini0852-chargeincoulombs, def:physics:phyx-mini-0852:phyxminiproblems-problemphyxmini0852-concentricconductingspheressetup, def:physics:phyx-mini-0852:phyxminiproblems-problemphyxmini0852-satisfiessphericalgausslaw}
  The conductor/Gauss-law model forces the inner-surface charge to cancel the central sphere's charge. This physically grounded intermediate conclusion contains no occurrence of the recorded answer
\end{lemma}
\begin{proof}
  The conclusion follows from the governing relations and figure readouts named in the statement of hollow inner surface charge cancels central.
\end{proof}
\begin{lemma}[central sphere charge is positive]
  \label{lem:physics:phyx-mini-0852:phyxminiproblems-problemphyxmini0852-central-sphere-charge-is-positive}
  \lean{PhyXMiniProblems.ProblemPhyXMini0852.central_sphere_charge_is_positive}
  \uses{def:physics:phyx-mini-0852:phyxminiproblems-problemphyxmini0852-chargeincoulombs, def:physics:phyx-mini-0852:phyxminiproblems-problemphyxmini0852-concentricconductingspheressetup, def:physics:phyx-mini-0852:phyxminiproblems-problemphyxmini0852-matchesprimaryconcentricspheresfigure, def:physics:phyx-mini-0852:phyxminiproblems-problemphyxmini0852-satisfiesconcentricsphericalgeometry, def:physics:phyx-mini-0852:phyxminiproblems-problemphyxmini0852-satisfiessphericalgausslaw, def:physics:phyx-mini-0852:phyxminiproblems-problemphyxmini0852-hasphysicalelectrostaticparameters}
  The nonzero outward cavity field, positive radius, and positive permittivity make the central sphere's charge positive. Together with the preceding lemma, this exposes the conflict between the primary figure and answer D
\end{lemma}
\begin{proof}
  The conclusion follows from the governing relations and figure readouts named in the statement of central sphere charge is positive.
\end{proof}
% --- Archon declaration topology end ---
% --- Archon physics formalization source end ---
